% ============================================================================
% Informe Técnico — Dron de superficie a vela sobre T-Beam ESP32
% Idioma: Español
% Autor (código original): Mohamed EL JILY
% Autor (refonte / firmware actual): Facundo Arito
% Compilación: pdflatex -interaction=nonstopmode informe_proyecto_es.tex (x2)
% ============================================================================

\documentclass[12pt,a4paper]{article}

% ── Codificación (pdflatex) ──────────────────────────────────────────────────
\usepackage[utf8]{inputenc}
\usepackage[T1]{fontenc}
\usepackage{lmodern}
\usepackage{textcomp}

% ── Idioma (sin babel-spanish para no depender de texlive-lang) ──────────────
\renewcommand{\contentsname}{\'Indice}
\renewcommand{\abstractname}{Resumen}
\renewcommand{\appendixname}{Anexo}
\renewcommand{\refname}{Referencias}
\renewcommand{\figurename}{Figura}
\renewcommand{\tablename}{Tabla}

% ── Diseño de página ─────────────────────────────────────────────────────────
\usepackage[top=2.5cm, bottom=2.5cm, left=2.8cm, right=2.8cm]{geometry}
\usepackage{parskip}
\usepackage{microtype}
\usepackage{setspace}
\onehalfspacing

% ── Colores ──────────────────────────────────────────────────────────────────
\usepackage[dvipsnames]{xcolor}
\definecolor{codebg}{HTML}{1E1E2E}
\definecolor{codefg}{HTML}{CDD6F4}
\definecolor{commentfg}{HTML}{6C7086}
\definecolor{keywordfg}{HTML}{CBA6F7}
\definecolor{stringfg}{HTML}{A6E3A1}
\definecolor{numberfg}{HTML}{FAB387}
\definecolor{warnbg}{HTML}{FFF8E1}
\definecolor{warnborder}{HTML}{F59E0B}
\definecolor{okbg}{HTML}{F0FDF4}
\definecolor{okborder}{HTML}{16A34A}
\definecolor{infobg}{HTML}{EFF6FF}
\definecolor{infoborder}{HTML}{2563EB}
\definecolor{tableheadbg}{HTML}{374151}
\definecolor{tableheadfg}{HTML}{FFFFFF}
\definecolor{tablerowalt}{HTML}{F9FAFB}

% ── Cajas de aviso ───────────────────────────────────────────────────────────
\usepackage{tcolorbox}
\tcbuselibrary{skins, breakable}

\newtcolorbox{warningbox}{
  colback=warnbg, colframe=warnborder,
  fonttitle=\bfseries, title={\color{warnborder}$\blacktriangle$\ Problema},
  breakable, boxrule=1pt, left=6pt, right=6pt
}
\newtcolorbox{okbox}{
  colback=okbg, colframe=okborder,
  fonttitle=\bfseries, title={\color{okborder}$\checkmark$\ Soluci\'on},
  breakable, boxrule=1pt, left=6pt, right=6pt
}
\newtcolorbox{infobox}[1][]{
  colback=infobg, colframe=infoborder,
  fonttitle=\bfseries, title={\color{infoborder}$\blacksquare$\ #1},
  breakable, boxrule=1pt, left=6pt, right=6pt
}

% ── Listings (código fuente) ─────────────────────────────────────────────────
\usepackage{listings}
\usepackage{lstautogobble}

\lstdefinestyle{cpp}{
  language=C++,
  backgroundcolor=\color{codebg},
  basicstyle=\ttfamily\footnotesize\color{codefg},
  keywordstyle=\color{keywordfg}\bfseries,
  stringstyle=\color{stringfg},
  commentstyle=\color{commentfg}\itshape,
  numberstyle=\tiny\color{numberfg},
  numbers=left, numbersep=8pt,
  breaklines=true, breakatwhitespace=true,
  frame=single, rulecolor=\color{codebg!40!black},
  tabsize=4, showstringspaces=false,
  autogobble=true,
  morekeywords={uint8_t,uint16_t,uint32_t,int8_t,int16_t,int32_t,bool,
                constexpr,noexcept,nullptr,override,IRAM_ATTR,portMUX_TYPE},
}
\lstdefinestyle{shell}{
  language=bash,
  backgroundcolor=\color{codebg},
  basicstyle=\ttfamily\footnotesize\color{codefg},
  breaklines=true, frame=single, rulecolor=\color{codebg!40!black},
  showstringspaces=false,
}
\lstdefinestyle{json}{
  backgroundcolor=\color{codebg},
  basicstyle=\ttfamily\footnotesize\color{codefg},
  breaklines=true, frame=single, rulecolor=\color{codebg!40!black},
  showstringspaces=false,
}

% ── Tablas ───────────────────────────────────────────────────────────────────
\usepackage{booktabs}
\usepackage{array}
\usepackage{longtable}
\usepackage{colortbl}
\usepackage{multirow}
\newcolumntype{L}[1]{>{\raggedright\arraybackslash}p{#1}}
\newcolumntype{C}[1]{>{\centering\arraybackslash}p{#1}}

% ── Hiperenlaces ─────────────────────────────────────────────────────────────
\usepackage[hidelinks]{hyperref}
\hypersetup{
  pdftitle={Informe T\'ecnico -- Dron de superficie a vela},
  pdfauthor={Facundo Arito},
  bookmarksnumbered=true,
}

% ── Formato de títulos ───────────────────────────────────────────────────────
\usepackage{titlesec}
\titleformat{\section}{\Large\bfseries\color{tableheadbg}}{\thesection.}{0.7em}{}[\vspace{-0.3em}\color{infoborder}\rule{\linewidth}{1.5pt}]
\titleformat{\subsection}{\large\bfseries\color{tableheadbg!80!black}}{\thesubsection.}{0.7em}{}
\titleformat{\subsubsection}{\normalsize\bfseries}{\thesubsubsection.}{0.6em}{}

% ── Listas ───────────────────────────────────────────────────────────────────
\usepackage{enumitem}
\setlist{itemsep=2pt, topsep=4pt}

% ── Encabezados y pies ───────────────────────────────────────────────────────
\usepackage{fancyhdr}
\setlength{\headheight}{15pt}
\pagestyle{fancy}
\fancyhf{}
\fancyhead[L]{\small Dron de superficie a vela --- Informe t\'ecnico}
\fancyhead[R]{\small \today}
\fancyfoot[C]{\small\thepage}
\renewcommand{\headrulewidth}{0.4pt}

% ── Utilidades ───────────────────────────────────────────────────────────────
\usepackage{graphicx}
\usepackage{amsmath}
\usepackage{amssymb}

\newcommand{\gpio}[1]{\texttt{GPIO#1}}
\newcommand{\us}[1]{\texttt{#1\,\textmu s}}
\newcommand{\file}[1]{\texttt{#1}}
\newcommand{\done}{\textcolor{okborder}{\textbf{[OK]}}}
\newcommand{\todo}{\textcolor{warnborder}{\textbf{[\ \ ]}}}

% ─────────────────────────────────────────────────────────────────────────────
\begin{document}

% ── Portada ──────────────────────────────────────────────────────────────────
\begin{titlepage}
  \centering
  \vspace*{2cm}
  {\Huge\bfseries Dron de superficie a vela\par}
  \vspace{0.5cm}
  {\Large Informe t\'ecnico --- Desarrollo del firmware y de la electr\'onica\par}
  \vspace{2cm}
  \rule{\linewidth}{1pt}
  \vspace{0.5cm}

  \begin{tabular}{ll}
    \textbf{Plataforma:} & LilyGO T-Beam V1.1 (ESP32 + SX1276 + GPS u-blox + AXP192) \\[4pt]
    \textbf{Lenguaje:}   & Arduino / C++, destino ESP32 Arduino Core 3.x \\[4pt]
    \textbf{Integrantes del grupo: } & Facundo Arito \\[4pt]
    \textbf{Fecha:}      & \today \\
  \end{tabular}

  \vspace{0.5cm}
  \rule{\linewidth}{1pt}
  \vspace{2cm}

  \begin{abstract}
    Este documento recorre la totalidad del desarrollo inform\'atico y electr\'onico
    del dron de superficie a vela basado en el m\'odulo T-Beam de Espressif/LilyGO.
    Cubre el concepto f\'isico del veh\'iculo, el c\'odigo original y sus limitaciones,
    la reescritura completa de la arquitectura de software, las dos revisiones del
    PCB, el modo manual unificado con propulsi\'on bidireccional, la navegaci\'on
    aut\'onoma a vela, la interfaz de estaci\'on terrena (IHM) y, de forma detallada,
    \textbf{todos los problemas} encontrados a lo largo del proyecto junto con las
    soluciones implementadas. Una secci\'on final reúne los problemas a\'un abiertos
    y las opciones de mejora propuestas para resolverlos.
  \end{abstract}

  \vfill
  \textit{Proyecto de investigaci\'on (PER) --- Dron aut\'onomo de adquisici\'on de datos marinos}
\end{titlepage}

\tableofcontents
\newpage

% =============================================================================
\section{Gu\'ia de uso para el pr\'oximo grupo de trabajo}
\label{sec:guia}
% =============================================================================

Esta secci\'on es una \textbf{gu\'ia pr\'actica de puesta en marcha} pensada para
el siguiente grupo que retome el proyecto. Explica qu\'e material hace falta, c\'omo
encender y pilotar el dron, c\'omo programar y conectar las dos placas T-Beam, y
c\'omo levantar la interfaz de estaci\'on terrena (IHM) tanto en Linux como en
Windows. Se recomienda leerla \textbf{completa} antes del primer encendido.

\subsection{Material necesario}

\begin{longtable}{L{4.5cm} L{8.5cm}}
  \toprule
  \rowcolor{tableheadbg}
  \textcolor{tableheadfg}{\textbf{Elemento}} &
  \textcolor{tableheadfg}{\textbf{Detalle}} \\
  \midrule
  \endhead
  \bottomrule
  \endfoot
  Dron (``barco'') & T-Beam montada en el casco sobre el PCB V2. Conversor
                     USB-serie CP2104, n\'umero de serie \texttt{02126CF1} \\[3pt]
  Transceptor (estaci\'on tierra) & Segunda T-Beam con el sketch
                     \file{transceiver.ino}. CP2104 serie \texttt{01C00B54} \\[3pt]
  Radio RC & Emisor Pro-Tronik PTR-6A + receptor R8X (2,4~GHz) \\[3pt]
  Bater\'ia & LiPo 2S \textbf{cargada} ($\approx$8,4~V) + portabater\'ias del
              casco. \textbf{Imprescindible} para un funcionamiento estable \\[3pt]
  Computadora & Linux (recomendado) o Windows, con Docker, Python 3 y navegador \\[3pt]
  Cables USB & Uno de buena calidad para el transceptor; otro para programar /
               alimentar el barco \\
\end{longtable}

\subsection{Puesta en marcha del dron (barco)}

\begin{enumerate}
  \item \textbf{Cargar la LiPo} por completo y conectarla al \textbf{terminal del PCB} del barco (esta indicado GND en la PCB).
  \item \textbf{Encender primero el emisor RC} y comprobar que el selector de modo
    (CH5) est\'a en \textbf{Manual} (posici\'on baja). El \textbf{primer ensayo debe
    ser siempre manual} para probar todos los actuadores.
  \item \textbf{Energizar el barco.} Conectar el USB de la T-Beam (ya sea a la alimentacion PCB o a una pc para leer el puerto serie) En el arranque imprime por serie l\'ineas de
    estado (p.\,ej.\ \texttt{[LORA] OK}).
  \item \textbf{Esperar el fix GPS:} el LED \texttt{TIMEPULSE} del NEO-6M parpadea a
    1~Hz cuando hay fix. Con LiPo conectada el fix llega en segundos (arranque en
    caliente); sin LiPo, cada encendido es un arranque en fr\'io de $\approx$2~min.
  \item \textbf{Pilotar en manual:}
  \begin{itemize}
    \item \textbf{CH5 (modo):} alto $=$ Autom\'atico, medio $=$ Sail (inerte, todo en
      neutro), bajo $=$ Manual.
    \item \textbf{CH2 (vela):} bascula la vela entre babor y estribor ($\pm10^{\circ}$).
    \item \textbf{CH4 (tim\'on):} posiciona el winch del tim\'on ($\pm90^{\circ}$).
    \item \textbf{CH3 (motor, bidireccional):} centro $=$ paro, un extremo $=$
      $+100~\%$ adelante, el otro $=$ $-100~\%$ atr\'as. Requiere el ESC en modo
      bidireccional (ver punto siguiente).
  \end{itemize}
\end{enumerate}

\begin{warningbox}
  \textbf{Seguridad:} si se pierde la se\~nal RC, el dron entra en \emph{Failsafe},
  que \textbf{navega de forma aut\'onoma}. Como el modo aut\'onomo nunca se valid\'o
  en el agua, el primer ensayo debe hacerse en manual, en aguas tranquilas y con un
  plan de recuperaci\'on (p.\,ej.\ bote o cabo).
\end{warningbox}

\subsection{Programaci\'on (flasheo) de las T-Beam}

Ambas placas usan el mismo n\'ucleo (\texttt{FQBN = esp32:esp32:t-beam}). El barco
lleva el sketch \file{main/}; el transceptor lleva \file{transceiver/}.

\begin{lstlisting}[style=shell, caption={Compilar y flashear con arduino-cli}]
# (rutas relativas a la carpeta raiz del proyecto = Post-Claude)
# Compilar el firmware del barco
~/bin/arduino-cli compile --fqbn esp32:esp32:t-beam main

# Flashear el barco (ajustar el puerto)
~/bin/arduino-cli upload --port /dev/ttyUSB0 --fqbn esp32:esp32:t-beam main

# Idem para el transceptor
~/bin/arduino-cli upload --port /dev/ttyUSB0 --fqbn esp32:esp32:t-beam transceiver
\end{lstlisting}

\begin{itemize}
  \item \textbf{Identificar cada placa por su n\'umero de serie} (en Linux:
    \texttt{ls -l /dev/serial/by-id/}). Barco $=$ \texttt{02126CF1},
    transceptor $=$ \texttt{01C00B54}.
  \item \textbf{Si el flasheo falla} (``\emph{chip stopped responding}'', el puerto
    se re-enumera): es un \emph{brownout} por alimentaci\'on insuficiente.  Entrar en modo bootloader manual (mantener \textbf{BOOT/IO0}
    mientras se pulsa \textbf{RST} o se reconecta el USB). Ver Secci\'on~\ref{sec:alim}.
  \item No tener un monitor serie abierto sobre el puerto mientras se programa.
\end{itemize}

\subsection{Conexi\'on de las dos T-Beam}

\begin{itemize}
  \item El \textbf{transceptor} se conecta por \textbf{USB a la computadora} de
    tierra. Aparece como \texttt{/dev/ttyUSBx} en Linux o \texttt{COMx} en Windows.
  \item El \textbf{barco no se conecta por cable} durante la operaci\'on: se comunica
    con el transceptor por \textbf{radio LoRa a 433~MHz}. Ambos firmwares ya usan los
    mismos par\'ametros de radio (433~MHz, SF7, BW125~kHz).
  \item El transceptor reenv\'ia a la IHM cada heartbeat del barco (y le inyecta el
    RSSI medido en tierra) y transmite al barco los comandos de la IHM.
\end{itemize}

\begin{warningbox}
  \textbf{No abrir} un monitor serie (gtkterm, el monitor del IDE Arduino, un script
  pyserial, etc.) sobre el puerto del transceptor \textbf{mientras corre la IHM}: los
  dos lectores compiten por el puerto y corrompen el JSON de telemetr\'ia.
\end{warningbox}

\subsection{Uso de la IHM en Linux}

\subsubsection{Programas que hay que instalar (Linux)}

\begin{longtable}{L{4.2cm} L{8.8cm}}
  \toprule
  \rowcolor{tableheadbg}
  \textcolor{tableheadfg}{\textbf{Programa}} &
  \textcolor{tableheadfg}{\textbf{Para qu\'e / c\'omo}} \\
  \midrule
  \endhead
  \bottomrule
  \endfoot
  \textbf{Docker Engine} + plugin \texttt{docker compose} & Levanta la base de datos
    \textbf{MongoDB} (imagen \texttt{mongo:7}, que se descarga sola la primera vez).
    En Ubuntu/Debian: \texttt{sudo apt install docker.io docker-compose-plugin}.
    A\~nadir el usuario al grupo: \texttt{sudo usermod -aG docker \$USER} (y reabrir
    sesi\'on) para no usar \texttt{sudo} \\[3pt]
  \textbf{Python 3} + \texttt{venv} + \texttt{pip} & Ejecuta el servidor y el puente
    serie. \texttt{sudo apt install python3 python3-venv python3-pip} \\[3pt]
  \textbf{Paquetes Python} & Se instalan dentro de un entorno virtual a partir de
    \file{requirements.txt}: \texttt{fastapi}, \texttt{uvicorn},
    \texttt{pymongo[asyncio]}, \texttt{pyserial}, \texttt{python-dotenv} \\[3pt]
  \textbf{Navegador web} & Para abrir la interfaz (Firefox, Chrome, etc.) \\
\end{longtable}

\textbf{Preparaci\'on (una sola vez)} --- crear el entorno virtual e instalar los
paquetes Python (rutas relativas a la ra\'iz \file{Post-Claude}):

\begin{lstlisting}[style=shell, caption={Instalacion inicial de la IHM en Linux}]
cd IHM
python3 -m venv .venv
.venv/bin/pip install -r requirements.txt
\end{lstlisting}

\textbf{Arranque y parada} --- conectar el transceptor por USB y, desde la carpeta
\file{IHM/}, ejecutar:

\begin{lstlisting}[style=shell, caption={Arrancar y detener la IHM en Linux}]
cd IHM
./start_ihm.sh                 # autodetecta el transceptor por su serie
# o forzar el puerto:
./start_ihm.sh /dev/ttyUSB0

# ... al terminar:
./stop_ihm.sh                  # libera el puerto 5000 y baja MongoDB (Docker)
\end{lstlisting}

\begin{itemize}
  \item \file{start\_ihm.sh} levanta MongoDB (contenedor Docker via
    \texttt{docker compose}), \file{serial\_link.py} (lee/escribe el puerto serie) y
    \file{webserver.py} (servidor web), y abre el navegador en
    \texttt{http://localhost:5000}.
  \item En el mapa: ver telemetr\'ia en vivo, definir waypoints con clic, enviar la
    ruta y los comandos (navegar, parar, medir viento, etc.).
  \item Conmutador \textbf{Sim / R\'eel}: en ``R\'eel'' usa el transceptor f\'isico;
    en ``Sim'' levanta un simulador de barco (sin hardware).
  \item El puerto del transceptor se \textbf{autodetecta} por su n\'umero de serie
    (\texttt{01C00B54}), por lo que normalmente no hace falta indicarlo.
\end{itemize}

\subsection{Uso de la IHM en Windows}

La IHM se desarroll\'o en Linux (los scripts de arranque son \texttt{bash}), por lo
que en Windows lo m\'as sencillo es ejecutarla dentro de WSL2.

\subsubsection{Programas que hay que instalar (Windows)}

\begin{longtable}{L{4.2cm} L{8.8cm}}
  \toprule
  \rowcolor{tableheadbg}
  \textcolor{tableheadfg}{\textbf{Programa}} &
  \textcolor{tableheadfg}{\textbf{Para qu\'e / c\'omo}} \\
  \midrule
  \endhead
  \bottomrule
  \endfoot
  \textbf{WSL2} (Ubuntu) & Subsistema Linux dentro de Windows; permite usar los
    mismos scripts \texttt{bash}. Instalar con \texttt{wsl --install} en PowerShell
    (como administrador) \\[3pt]
  \textbf{Docker Desktop} & Levanta MongoDB. Instalarlo y activar la
    \textbf{integraci\'on con WSL2} en sus ajustes \\[3pt]
  \textbf{usbipd-win} & Comparte el puerto USB del transceptor con WSL (Windows no lo
    pasa solo). Instalar con \texttt{winget install usbipd} \\[3pt]
  \textbf{Python 3} + paquetes & Ya viene en Ubuntu de WSL; instalar los paquetes
    igual que en Linux (\file{requirements.txt}). Si se usa Windows nativo, instalar
    Python desde \texttt{python.org} \\[3pt]
  \textbf{Navegador web} & Para abrir \texttt{http://localhost:5000} \\
\end{longtable}

\textbf{Camino recomendado --- WSL2 + Docker Desktop:} dentro de Ubuntu (WSL) se
siguen \textbf{exactamente los mismos pasos que en Linux} (crear el \texttt{.venv},
instalar \file{requirements.txt}, \file{./start\_ihm.sh}). Para que WSL ``vea'' el
transceptor hay que adjuntar el USB desde PowerShell:

\begin{lstlisting}[style=shell, caption={Adjuntar el USB del transceptor a WSL (PowerShell admin)}]
usbipd list                       # ver el BUSID del CP210x (transceptor)
usbipd bind --busid <BUSID>
usbipd attach --wsl --busid <BUSID>   # dentro de WSL aparece como /dev/ttyUSBx
\end{lstlisting}

\textbf{Alternativa --- Windows nativo (sin WSL):} instalar Docker Desktop y Python
para Windows y lanzar los componentes a mano: \texttt{docker compose up} para
MongoDB, luego \texttt{python serial\_link.py} y \texttt{python webserver.py}. El
puerto serie ser\'a \texttt{COMx} (verificarlo en el Administrador de dispositivos)
y hay que ajustar \texttt{SERIAL\_PORT\_REAL} en el archivo \file{.env} a ese
\texttt{COMx}. En ambos casos la interfaz se abre en \texttt{http://localhost:5000}.

\begin{warningbox}
  \textbf{Red WiFi --- EDUROAM no funciona.} Como la IHM usa MongoDB (en Docker) y la
  pila de red local, en la red institucional \textbf{EDUROAM la IHM no funciona}:
  EDUROAM aplica aislamiento de clientes y restricciones de red (firewall/proxy) que
  impiden el correcto funcionamiento de Docker/MongoDB y de la comunicaci\'on local
  necesaria.

  \textbf{Soluci\'on (probada por el grupo):} conectar la computadora que ejecuta la
  IHM a una red sin esas restricciones. En la pr\'actica, \textbf{compartir los datos
  m\'oviles del celular (hotspot / anclaje a red) a la computadora}: con el hotspot
  del tel\'efono la IHM levanta y funciona correctamente.
\end{warningbox}

% =============================================================================
\section{Presentaci\'on del proyecto}
% =============================================================================

\subsection{Concepto físico}

El dron es un veh\'iculo de superficie propulsado principalmente por una vela.
Su concepci\'on es original: la vela es de \textbf{perfil de ala} y gira libremente
alrededor de un eje vertical solidario al casco. Su aerodin\'amica pasiva orienta
autom\'aticamente el perfil de modo que el viento incide siempre a unos
$45^{\circ}$ del eje longitudinal del dron, sin necesidad de un control activo del
\'angulo de la vela. La maniobra se realiza con dos actuadores principales:

\begin{itemize}
  \item \textbf{Servo aler\'on (vela) --- Futaba S3003:} controla un peque\~no
    flap en el borde de fuga de la vela. Su recorrido \'util es estrictamente de
    $\pm10^{\circ}$. La deflexi\'on determina hacia qu\'e lado (babor o estribor)
    empuja la vela, eligiendo de hecho la amura sobre la que navega el barco.
    \textbf{No es un control continuo}: solo las dos posiciones extremas tienen
    sentido f\'isico, por lo que se modela como un estado binario ($+1$ / $-1$).

  \item \textbf{Servo tim\'on / rotor --- Graupner Regatta ECO II (mod.\ 5176):}
    es el actuador principal de rumbo (guini\'ada). Es un \emph{winch de vela
    multivuelta posicional}: la se\~nal PWM ordena una \textbf{posici\'on} del eje
    (no una velocidad). Centro mec\'anico = \us{1500}, recorrido total $=$ 6 vueltas,
    y mantiene la posici\'on bajo carga como un servo est\'andar.

  \item \textbf{Propulsor (ESC):} utilizado \'unicamente cuando no hay viento
    suficiente. Permite la propulsi\'on a h\'elice como respaldo. En el PCB V2 se
    usa un \'unico ESC (se elimin\'o el segundo).
\end{itemize}

\begin{infobox}[Principio de navegaci\'on sin sensor de viento]
  El sistema \textbf{no mide} la direcci\'on del viento directamente: la
  \emph{deduce} a partir del rumbo GPS (track) y del estado de los actuadores.
  La estimaci\'on requiere un desplazamiento m\'inimo de 30~m. Esta carencia de
  sensores es central en el proyecto y se retoma en la Secci\'on~\ref{sec:abiertos}.
\end{infobox}

\subsection{Objetivo del proyecto}

El proyecto es un PER (Proyecto de Estudios y Recherche) cuyo prop\'osito final es
la \textbf{adquisici\'on aut\'onoma de datos marinos}: el dron debe alcanzar puntos
de paso (\emph{waypoints}) GPS navegando a vela, recolectar datos de sensores y
transmitirlos por radio LoRa a una estaci\'on terrena.

\subsection{Material principal}

\begin{longtable}{L{4cm} L{9cm}}
  \toprule
  \rowcolor{tableheadbg}
  \textcolor{tableheadfg}{\textbf{Componente}} &
  \textcolor{tableheadfg}{\textbf{Descripci\'on}} \\
  \midrule
  \endhead
  \bottomrule
  \endfoot
  LilyGO T-Beam V1.1 & M\'odulo que integra ESP32 WROVER, LoRa SX1276 433~MHz,
                       GPS u-blox NEO-6M, AXP192 (gesti\'on de energ\'ia) y
                       portabater\'ias 18650 \\[3pt]
  Futaba S3003 & Servo aler\'on de la vela, centro a \us{1520}, $\pm10^{\circ}$ \\[3pt]
  Graupner Regatta ECO II & Winch de tim\'on multivuelta, \us{1500} = centro,
                            1000--2000~\textmu s = rango, hasta 3,3~A en carga \\[3pt]
  Pro-Tronik PTR-6A & Emisor 6 v\'ias FHSS 2,4~GHz, PWM 1000--2000~\textmu s \\[3pt]
  Pro-Tronik R8X & Receptor 8 v\'ias, PWM est\'andar 50~Hz \\[3pt]
  Pro-Tronik Black Fet & ESC de la serie, programaci\'on por tarjeta EPRG-3 \\[3pt]
  Antena GPS & Taoglas activa, conector u.FL \\[3pt]
  PCB v1 / PCB v2 & Placa de interfaz a medida (KiCad 9) \\
\end{longtable}

% =============================================================================
\section{C\'odigo original --- \texttt{boat.ino}}
% =============================================================================

\subsection{Arquitectura y funcionamiento}

El c\'odigo original (\file{boat.ino}, del primer grupo, fuera de este
repositorio), escrito por Mohamed EL JILY, constituye el punto de partida del
proyecto. Estaba organizado
en m\'odulos: \file{Gps}, \file{ServoControl}, \file{RadioReceiver},
\file{MotorControl} y la l\'ogica de navegaci\'on en \file{boat.ino}. Inclu\'ia ya
una l\'ogica de navegaci\'on a vela (\'angulo relativo viento/rumbo, cambio de
amura autom\'atico, orzar/arribar) y recepci\'on de waypoints por LoRa en JSON.

\subsection{Problemas y limitaciones identificados}

Al retomar el proyecto se identificaron varios problemas bloqueantes que, en conjunto con la nesesidad de nuevas features discutidas con el cliente y el profesor, motivaron la reescritura completa del firmware.

\subsubsection{Biblioteca ESP32Servo --- conflicto LEDC / interrupciones}

\begin{warningbox}
  La biblioteca \texttt{ESP32Servo} utiliza el perif\'erico \textbf{LEDC} del ESP32.
  LEDC es \textbf{incompatible} con la decodificaci\'on de PWM por interrupciones
  GPIO usada para leer la radio: cuando LEDC est\'a activo, las interrupciones del
  receptor RC se perturban y los canales le\'idos \textbf{vuelven constantemente a
  cero}, aunque las se\~nales de entrada sean correctas en el oscilloscopio. Esto
  hace imposible controlar los servos y leer la radio al mismo tiempo. Fue el
  problema m\'as cr\'itico del proyecto.
\end{warningbox}

\subsubsection{Asignaci\'on de pines RC --- conflicto I2C / AXP192}

\begin{warningbox}
  El c\'odigo original asignaba canales RC a \gpio{21} y \gpio{22}, que son los
  pines I2C internos del T-Beam usados por el AXP192 (gestor de energ\'ia).
  Adem\'as, \gpio{23} est\'a conectado internamente al RESET del SX1276 (LoRa):
  conectar una interrupci\'on en ese pin pod\'ia provocar resets esp\'ureos de la
  radio.
\end{warningbox}

\subsubsection{Pin de bater\'ia \gpio{35} --- drenaje de 20~mA continuo}

El T-Beam V1.1 tiene un divisor resistivo de \emph{baja impedancia}
($\approx370~\Omega$ total) soldado de forma permanente en \gpio{35}. Aun sin leer
el ADC, ese divisor consume $\approx20$~mA continuos ($\approx1{,}8$~W) directamente
de la bater\'ia. \textbf{Nunca debe leerse \gpio{35}.}

\subsubsection{Pin LoRa RST --- error de documentaci\'on}

La documentaci\'on de LilyGO y muchos ejemplos indican \gpio{14} para el RESET del
SX1276. El hardware del T-Beam V1.1 conecta f\'isicamente el RESET en \gpio{23}.

\subsubsection{Gesti\'on de ESC --- sin armado ni limitaci\'on de variaci\'on}

El c\'odigo original no implementaba ning\'una funcion del ESC dado que fue una feature a futuro, al igual que la medicion de bateria.

% =============================================================================
\section{PCB Versi\'on 1 --- Placa de interfaz original}
% =============================================================================

La primera placa realizaba la interfaz entre el T-Beam y los actuadores
(conectores de 3 pines para servos/ESC, conectores RC, regulador reductor 5~V
Pololu, bornes de potencia). Sus problemas se resumen a continuaci\'on.

\begin{longtable}{L{3.5cm} L{9.5cm}}
  \toprule
  \rowcolor{tableheadbg}
  \textcolor{tableheadfg}{\textbf{Problema}} &
  \textcolor{tableheadfg}{\textbf{Impacto}} \\
  \midrule
  \endhead
  \bottomrule
  \endfoot
  PWM 3,3~V sin amplificaci\'on de nivel & Servos y ESC (l\'ogica 5~V) pueden no
    detectar correctamente los pulsos \\[3pt]
  Pines RC en \gpio{21}/\gpio{22} & Conflicto I2C/AXP192 \\[3pt]
  ESC en \gpio{32}/\gpio{33} & Coinciden con LoRa DIO1/DIO2 del T-Beam \\[3pt]
  \gpio{35} para bater\'ia & Divisor de 370~$\Omega$ = 20~mA continuos \\[3pt]
  Sin protecci\'on de polaridad & Riesgo ante alimentaci\'on incorrecta \\
\end{longtable}

% =============================================================================
\section{Reescritura del firmware --- Arquitectura Post-Claude}
% =============================================================================

\subsection{Motivaciones}

La reescritura completa busca corregir todas las limitaciones identificadas y
preparar la arquitectura para las fases aut\'onomas. Objetivos:

\begin{enumerate}
  \item Usar \textbf{MCPWM} en lugar de LEDC para las salidas de actuadores.
  \item Decodificaci\'on RC \textbf{enteramente por interrupciones} con protecci\'on ISR.
  \item Arquitectura orientada a objetos en capas, extensible por fases.
  \item Armado de ESC y limitaci\'on de variaci\'on (\emph{slew rate}).
  \item Gesti\'on limpia de la energ\'ia v\'ia AXP192.
  \item GPS, LoRa y bater\'ia funcionando simult\'aneamente.
\end{enumerate}

\subsection{Arquitectura en capas}

\begin{lstlisting}[style=shell, caption={\'Arbol del proyecto (raiz = Post-Claude)}]
Post-Claude/                   <- carpeta raiz del proyecto
  main/main.ino                <- Punto de entrada (begin / update)
  main/src/
    app/DroneApp.cpp           <- Orquestador (scheduler por software)
    config/BoardConfig.h       <- Asignacion de GPIO y umbrales RC
    config/Calibration.h       <- Valores en microsegundos (servo/ESC)
    core/Types.h               <- Tipos compartidos (RcFrame, ActuatorCommand)
    control/ModeManager        <- CH5 -> modo de control
    control/ManualController   <- Modo manual unificado (vela+timon+motor)
    control/AutoController     <- Envoltura de la navegacion autonoma
    drivers/RcReceiver         <- Lectura PWM por interrupciones
    drivers/McpwmActuators     <- Salidas MCPWM (vela, timon, ESC)
    drivers/AxpPower           <- Init AXP192 (rieles GPS/LoRa/3.3V)
    drivers/BatteryAdc         <- ADC de tension de bateria
    drivers/GpsUart            <- Lectura GPS (TinyGPSPlus + UBX)
    drivers/LoRaRadio          <- Driver SX1276
    comm/LoRaComm              <- Protocolo JSON LoRa (heartbeat + comandos)
    navigation/                <- navigation.h, Navigator, MissionManager...
  transceiver/transceiver.ino  <- Estacion terrena (puente USB <-> LoRa)
  IHM/                         <- Interfaz web de estacion terrena (FastAPI)
\end{lstlisting}

El reparto sigue una l\'ogica clara: \emph{drivers} (hardware), \emph{servicios /
comm} (datos), \emph{navigation} (geometr\'ia y misi\'on), \emph{control}
(decisiones) y \emph{app} (orquestaci\'on). El planificador es cooperativo (sin
RTOS): \file{DroneApp::update()} dispara cada tarea cuando su per\'iodo vence,
evitando \texttt{delay()} que bloquear\'ia la lectura RF.

\subsection{Decisiones t\'ecnicas fundamentales}

\subsubsection{MCPWM obligatorio --- nunca LEDC}

\begin{lstlisting}[style=cpp, caption={McpwmActuators.cpp -- inicializacion MCPWM}]
bool McpwmActuators::begin() {
    mcpwm_gpio_init(MCPWM_UNIT_0, MCPWM0A, BoardConfig::SAIL_SERVO_PIN);
    mcpwm_gpio_init(MCPWM_UNIT_0, MCPWM0B, BoardConfig::ROTOR_SERVO_PIN);
    mcpwm_gpio_init(MCPWM_UNIT_0, MCPWM1A, BoardConfig::ESC1_PIN);
    mcpwm_config_t cfg = {};
    cfg.frequency    = Calibration::PWM_FREQ_HZ;   // 50 Hz
    cfg.duty_mode    = MCPWM_DUTY_MODE_0;          // activo alto
    cfg.counter_mode = MCPWM_UP_COUNTER;
    mcpwm_init(MCPWM_UNIT_0, MCPWM_TIMER_0, &cfg);  // vela + timon
    mcpwm_init(MCPWM_UNIT_0, MCPWM_TIMER_1, &cfg);  // ESC
    ...
}
\end{lstlisting}

El perif\'erico MCPWM es independiente de las interrupciones GPIO: ambos coexisten
sin perturbarse. Esta restricci\'on es \textbf{innegociable} en este proyecto.

\subsubsection{Interrupciones RC con protecci\'on portMUX}

\begin{lstlisting}[style=cpp, caption={RcReceiver.cpp -- ISR y seccion critica}]
void IRAM_ATTR RcReceiver::handleEdge(uint8_t pin, Ch ch) {
    const uint32_t now = micros();
    const int      idx = static_cast<int>(ch);
    if (gpio_get_level(static_cast<gpio_num_t>(pin))) {
        portENTER_CRITICAL_ISR(&mux_);
        riseUs_[idx] = now;                        // flanco de subida
        portEXIT_CRITICAL_ISR(&mux_);
    } else {
        uint32_t rise;
        portENTER_CRITICAL_ISR(&mux_);
        rise = riseUs_[idx];
        portEXIT_CRITICAL_ISR(&mux_);
        const uint32_t width = now - rise;         // ancho de pulso
        if (width >= RC_VALID_MIN_US && width <= RC_VALID_MAX_US) {
            portENTER_CRITICAL_ISR(&mux_);
            pulseUs_[idx]     = static_cast<uint16_t>(width);
            lastValidUs_[idx] = now;
            portEXIT_CRITICAL_ISR(&mux_);
        }
    }
}
\end{lstlisting}

Las ISR se ubican en RAM (\texttt{IRAM\_ATTR}) y se sincronizan con la
boucle principal mediante \texttt{portMUX} (mecanismo FreeRTOS apropiado para el
ESP32 dual-core). Un \emph{timeout} de 100~ms marca un canal como perdido.

\subsubsection{Todas las consignas de actuadores en microsegundos}

Todos los actuadores se comandan exclusivamente en microsegundos, nunca en grados.
Esto elimina una capa de conversi\'on propensa a errores y corresponde
directamente al protocolo PWM de servos y ESC.

\subsubsection{Slew rate --- protecci\'on contra tirones}

\begin{lstlisting}[style=cpp, caption={McpwmActuators.cpp -- limitacion de variacion}]
uint16_t McpwmActuators::slew(uint16_t current, uint16_t target, uint16_t step) {
    if (current < target) {
        uint32_t next = (uint32_t)current + step;
        return (next >= target) ? target : (uint16_t)next;
    }
    if (current > target)
        return (current > target + step) ? (uint16_t)(current - step) : target;
    return current;
}
// ESC_SLEW_US = 30 us/tick, llamado cada 20 ms -> max 1500 us/s
\end{lstlisting}

% =============================================================================
\section{Fase 1 --- Control manual por radio}
% =============================================================================

\subsection{Realizaci\'on}

La Fase~1 cubre el control manual completo. Se implementaron los drivers
(\file{AxpPower}, \file{RcReceiver}, \file{McpwmActuators}), el \file{ModeManager},
el \file{ManualController} y el orquestador \file{DroneApp} (ticks de 20~ms para
control, 200~ms para debug, 1~s para LoRa).

\subsection{Modos de vuelo (selector CH5) --- versi\'on unificada actual}

El selector de 3 posiciones est\'a cableado en CH5. La distribuci\'on de modos fue
\textbf{reorganizada} respecto a la versi\'on inicial (que ten\'ia dos modos
manuales separados ``ManualServo'' y ``ManualProp''):

\begin{longtable}{C{2.6cm} C{2.4cm} L{8cm}}
  \toprule
  \rowcolor{tableheadbg}
  \textcolor{tableheadfg}{\textbf{CH5}} &
  \textcolor{tableheadfg}{\textbf{Modo}} &
  \textcolor{tableheadfg}{\textbf{Comportamiento}} \\
  \midrule
  \endhead
  \bottomrule
  \endfoot
  $\leq$ \us{1250} & Automatic & Navegaci\'on aut\'onoma (sigue la misi\'on LoRa) \\[3pt]
  1400--1600 & Sail & \textbf{Inerte}: todas las salidas en neutro \\[3pt]
  $>$ \us{1800} & Manual & Modo manual \textbf{unificado} (ver abajo) \\[3pt]
  Se\~nal perdida & Failsafe & Sigue la misi\'on aut\'onoma (como Automatic) \\
\end{longtable}

\begin{warningbox}
  \textbf{Atenci\'on de seguridad:} en p\'erdida de RC el sistema entra en
  \emph{Failsafe}, que se comporta como el modo aut\'onomo (navega solo). Como el
  modo aut\'onomo nunca se prob\'o en el agua, el \textbf{primer ensayo real debe
  ser manual} y con plan de recuperaci\'on.
\end{warningbox}

\subsection{Modo manual unificado}

Un \'unico modo (\texttt{computeManual}) controla a la vez la vela, el tim\'on y el
propulsor:

\begin{itemize}
  \item \textbf{CH2 $\to$ vela (binaria):} banda muerta $\pm35$~\textmu s alrededor
    de \us{1500}. El primer frame inicializa el estado desde la posici\'on del
    manche para evitar un salto al entrar en el modo. Dentro de la banda muerta se
    mantiene el \'ultimo estado.
  \item \textbf{CH4 $\to$ tim\'on (posicional):} el recorrido medido CH4
    (1180--1790~\textmu s) se mapea al rango f\'isico del winch
    (\texttt{ROTOR\_MIN/MAX\_US} = 1417--1583~\textmu s, $\pm90^{\circ}$). Banda muerta
    central que fija el centro exacto (\us{1500}).
  \item \textbf{CH3 $\to$ propulsor (bidireccional, ver Secci\'on~\ref{sec:bidir}).}
\end{itemize}

\begin{lstlisting}[style=cpp, caption={ManualController -- logica de vela binaria con estado memorizado}]
if (!sailStateInitialized_) {
    sailState_ = (frame.ch2 >= RC_MID_US) ? +1 : -1;
    sailStateInitialized_ = true;
}
if (frame.ch2 > RC_MID_US + db)      sailState_ = +1;
else if (frame.ch2 < RC_MID_US - db) sailState_ = -1;
cmd.sailUs = (sailState_ > 0) ? SAIL_PLUS_US : SAIL_MINUS_US;
\end{lstlisting}

\subsection{Propulsor bidireccional (marcha adelante / atr\'as)}
\label{sec:bidir}

En la versi\'on inicial el motor solo iba de 0 a 100~\% (control inverso, con el
ESC en modo unidireccional, \us{1000} = paro). Se modific\'o para que el motor vaya
de \textbf{$+100~\%$ (adelante) a $-100~\%$ (atr\'as)} usando el ESC en modo
\textbf{bidireccional}, en el que el neutro pasa a \us{1500}:

\begin{longtable}{C{4.5cm} C{3cm} L{5cm}}
  \toprule
  \rowcolor{tableheadbg}
  \textcolor{tableheadfg}{\textbf{Posici\'on del manche CH3}} &
  \textcolor{tableheadfg}{\textbf{Potencia}} &
  \textcolor{tableheadfg}{\textbf{ESC}} \\
  \midrule
  \endhead
  \bottomrule
  \endfoot
  \texttt{CH3\_FULL\_US} (\us{1100}) & $+100~\%$ adelante & \us{2000} \\[3pt]
  \texttt{CH3\_CENTER\_US} (\us{1545}) $\pm$ banda muerta & 0~\% paro & \us{1500} \\[3pt]
  \texttt{CH3\_ZERO\_US} (\us{1990}) & $-100~\%$ atr\'as & \us{1000} \\
\end{longtable}

Como el manche de gas es de tipo \emph{trinquete} (no se autocentra), el paro se
ubica en el centro de su recorrido con una banda muerta
(\texttt{ESC\_CENTER\_DEADBAND\_US} $=40$~\textmu s). Para mantener coherente todo el
firmware, el neutro \us{1500} se usa tambi\'en en el arranque, en Failsafe, en el
modo Sail y como paro del modo aut\'onomo; el l\'imite inferior de saturaci\'on se
baj\'o a \us{1000} para permitir la marcha atr\'as.

\begin{warningbox}
  La marcha atr\'as solo funciona f\'isicamente si el ESC est\'a \textbf{programado
  en modo bidireccional} con la tarjeta EPRG-3. Si no, todo valor por debajo de
  \us{1500} se ignora y solo se obtiene marcha adelante.
\end{warningbox}

\subsection{Resultados de las pruebas --- Fase 1}

\begin{itemize}
  \item \done\ Aler\'on conmuta correctamente entre \us{1465} y \us{1575}.
  \item \done\ Tim\'on responde proporcionalmente al manche CH4.
  \item \done\ Failsafe se dispara en menos de 100~ms tras la p\'erdida de se\~nal.
  \item \done\ Compilaci\'on \texttt{arduino-cli}: $\approx$28~\% flash, 7~\% RAM.
  \item \todo\ Calibraci\'on de banco de \texttt{SAIL\_PLUS/MINUS\_US} y de
    \texttt{CH3\_CENTER\_US} pendiente.
\end{itemize}

% =============================================================================
\section{Fase 2 --- GPS y navegaci\'on}
% =============================================================================

Se implementaron \file{GpsUart} (Serial1 en \gpio{34}/\gpio{12}, TinyGPSPlus),
\file{Navigator} (haversine: distancia y rumbo), \file{MissionPlan} (hasta 16
waypoints) y \file{MissionManager} (m\'aquina de estados
Idle$\to$Running$\to$Returning$\to$Complete). Dos problemas cr\'iticos de GPS se
resolvieron en esta fase.

\subsection{Problema cr\'itico --- GPS mudo tras configuraci\'on en flash}

\begin{warningbox}
  \textbf{S\'intoma:} el contador \texttt{chars} se queda bloqueado, el GPS est\'a
  alimentado pero no produce ninguna frase NMEA.

  \textbf{Causa:} una sesi\'on previa (IDE Arduino o u-center) guard\'o en la flash
  interna del NEO-6M una configuraci\'on CFG-PRT que desactiva la salida NMEA. Esa
  configuraci\'on sobrevive a los cortes de alimentaci\'on.
\end{warningbox}

\begin{okbox}
  Enviar un reset de f\'abrica \texttt{UBX CFG-CFG} al arrancar (borra la
  configuraci\'on de flash y carga los valores de ROM). Implementado en
  \file{GpsUart::begin()}.
\end{okbox}

\subsection{Problema cr\'itico --- antena activa sin alimentaci\'on}
\label{subsec:antena}

\begin{warningbox}
  \textbf{S\'intoma:} NMEA recibido pero \texttt{visible = 0} sat\'elites incluso
  en exterior despejado y con antena externa conectada.

  \textbf{Causa:} el supervisor de antena del NEO-6M est\'a desactivado por
  defecto. Sin activarlo, no se entrega tensi\'on de polarizaci\'on al LNA de la
  antena activa y el conmutador RF del T-Beam permanece en la antena cer\'amica
  interna, cortocircuitando el conector u.FL.
\end{warningbox}

\begin{okbox}
  Enviar \texttt{UBX CFG-ANT} con \texttt{flags = 0x001B} (control de tensi\'on RF,
  detecci\'on de cortocircuito, recuperaci\'on). Tras la correcci\'on: fix en menos
  de 60~s, 7 sat\'elites, HDOP $\approx$2,2.
\end{okbox}

\subsection{Nota sobre el arranque en fr\'io}

Sin bater\'ia LiPo conectada, cada encendido pierde el almanaque GPS: el arranque
en fr\'io tarda $\approx$2~min en exterior. Con bater\'ia se conservan las
ef\'emerides y el fix se obtiene en segundos (arranque en caliente).

% =============================================================================
\section{Fase 3 --- LoRa y telemetr\'ia}
% =============================================================================

Se implement\'o el driver \file{LoRaRadio} (SX1276, SPI HSPI, 433~MHz), la capa de
protocolo \file{LoRaComm} (heartbeat JSON a 1~Hz + despachador de comandos) y el
sketch de estaci\'on terrena \file{transceiver.ino} (puente USB$\leftrightarrow$LoRa
con CLI compacta). Todo verificado en hardware (RSSI $-101$ a $-106$~dBm a corta
distancia).

\subsection{Protocolo JSON --- heartbeat actual}

El formato se afin\'o para mantenerse holgadamente por debajo del l\'imite de
255~bytes de la FIFO del SX1276 (se eliminaron campos redundantes, se a\~nadieron
indicadores de salud GPS/RC, y el RSSI lo inyecta el transceptor en tierra):

\begin{lstlisting}[style=json, caption={Heartbeat dron -> tierra (1 Hz)}]
{"origin":"boat","type":"info","message":{
  "mode":"standby|route-ready|navigate",
  "location":[48.36000,-4.56000],
  "servos":{"sail":10,"rudder":-12},
  "heading":270,"wind":180,"bat":7.42,
  "fix":1,"sat":7,"hdop":2.2,"rc":1,
  "wt":3,"wc":1,"wobs":42,"rssi":-104}}
\end{lstlisting}

\begin{itemize}
  \item \texttt{servos.sail} $=\pm10^{\circ}$ (binaria); \texttt{servos.rudder} en
    grados del winch. \texttt{fix/sat/hdop} = salud GPS; \texttt{rc}=1 si el
    receptor entrega pulsos.
  \item \texttt{wt/wc} = waypoints total/actual. \texttt{wobs} = progreso de la
    medici\'on de viento (0--100~\%), presente \emph{solo} mientras se mide.
  \item \texttt{rssi} \textbf{no} lo env\'ia el barco: lo inyecta el transceptor en
    tierra, por lo que no cuenta en el presupuesto de 255~bytes.
\end{itemize}

\subsection{Comandos (tierra $\to$ dron)}

\begin{lstlisting}[style=json, caption={Comandos soportados}]
{"origin":"server","type":"command","message":{"navigate":true}}
{"origin":"server","type":"command","message":{"stop":true}}
{"origin":"server","type":"command","message":{"restart":true}}
{"origin":"server","type":"command","message":{"wind-observation":true}}
{"origin":"server","type":"command","message":{"wind-command":{"value":225}}}
{"origin":"server","type":"command","message":{"home":{"lat":48.36,"lon":-4.56}}}
{"origin":"server","type":"command","message":
  {"waypoints":{"number":2,"points":"48.36,-4.56,48.37,-4.55"}}}
\end{lstlisting}

\begin{infobox}[Nota sobre el TX bloqueante]
  \texttt{LoRa.endPacket()} es s\'incrono. A SF7/BW125~kHz, un paquete de
  $\approx$210~bytes tarda $\approx$450~ms, bloqueando la boucle principal una vez
  por segundo. Aceptable por ahora; conviene pasar a TX as\'incrono antes de la
  navegaci\'on aut\'onoma en agua si el jitter del lazo de control se vuelve
  cr\'itico.
\end{infobox}

% =============================================================================
\section{PCB Versi\'on 2 --- Mejoras de hardware}
% =============================================================================

La V2 del PCB corrige todos los problemas de la V1 e incorpora los aprendizajes
del desarrollo de software.

\subsection{Correspondencia de pines (PCB V2)}

\begin{longtable}{C{3cm} C{1.6cm} L{7.5cm}}
  \toprule
  \rowcolor{tableheadbg}
  \textcolor{tableheadfg}{\textbf{Se\~nal}} &
  \textcolor{tableheadfg}{\textbf{GPIO}} &
  \textcolor{tableheadfg}{\textbf{Notas}} \\
  \midrule
  \endhead
  \bottomrule
  \endfoot
  RC CH2 (aler\'on) & 39 & SVN, solo entrada \\[2pt]
  RC CH3 (gas) & 14 & Liberado (ya no comparte I2C) \\[2pt]
  RC CH4 (tim\'on) & 13 & Liberado \\[2pt]
  RC CH5 (modo) & 4 & Conmutador 3 posiciones \\[2pt]
  Servo vela (PWM1) & 2 & V\'ia transistor BC548 \\[2pt]
  Servo tim\'on (PWM2) & 25 & V\'ia transistor BC548 \\[2pt]
  ESC1 & 15 & V\'ia transistor BC548 (\'unico ESC) \\[2pt]
  ADC bater\'ia & 36 & SVP, ADC1\_CH0, solo entrada \\[2pt]
  LoRa SCK/MISO/MOSI/CS & 5/19/27/18 & SPI HSPI \\[2pt]
  LoRa RST & 23 & Ahora libre (CH6 eliminado) \\[2pt]
  LoRa IRQ (DIO0) & 26 & \\[2pt]
  GPS RX / TX & 34 / 12 & Serial1 \\[2pt]
  I2C SDA / SCL & 21 / 22 & AXP192 interno + \textbf{conector de expansi\'on} \\
\end{longtable}

\subsection{Amplificaci\'on de nivel l\'ogico --- transistores BC548}
\label{subsec:bc548}

Las salidas PWM del ESP32 son de 3,3~V; los servos esperan 5~V. La V2 usa
transistores NPN BC548 en emisor com\'un.

\begin{infobox}[Descubrimiento importante --- DUTY\_MODE\_0 en ambos canales]
  Inicialmente el canal del tim\'on (OPR\_B) usaba \texttt{MCPWM\_DUTY\_MODE\_1}
  (inversi\'on), suponiendo que el NPN invert\'ia la se\~nal y hab\'ia que
  preinvertir. En realidad el montaje BC548 ya invierte; aplicar
  \texttt{DUTY\_MODE\_1} produc\'ia una \textbf{doble inversi\'on} que dejaba el
  servo del tim\'on inoperante. La soluci\'on es usar \texttt{DUTY\_MODE\_0}
  (activo alto) en \textbf{ambos} canales.
\end{infobox}

\subsection{Medici\'on de bater\'ia --- divisor de alta impedancia}

Se usa un divisor R5 $=562$~k$\Omega$ / R6 $=120$~k$\Omega$ (raz\'on $5{,}683$) en
\gpio{36} con atenuaci\'on de 11~dB. A 7,4~V el divisor consume
$\approx10{,}9$~\textmu A, despreciable.

\begin{warningbox}
  \textbf{Detalle de montaje:} en una placa las resistencias se soldaron
  invertidas (120~k arriba, 562~k abajo), dando una lectura de 17,9~V en lugar de
  $\approx$5~V. El c\'odigo era correcto; la soluci\'on fue invertir f\'isicamente
  las resistencias.
\end{warningbox}

\subsection{AXP192 --- I2C persistente}

En la V1 se llamaba \texttt{Wire.end()} tras inicializar el AXP192 para liberar
\gpio{21}/\gpio{22}. En la V2 esos pines quedan libres para I2C, por lo que
\texttt{Wire.end()} se elimina: el bus I2C permanece activo y disponible para
futuros sensores en el conector de expansi\'on.

% =============================================================================
\section{Navegaci\'on aut\'onoma (Fases 5 y 6)}
% =============================================================================

\subsection{Fase 5 --- Estimaci\'on de viento desde el track GPS}

Sin sensor de viento, la direcci\'on se infiere navegando sobre una amura fija:
el barco fija la vela en $+10^{\circ}$, suaviza el rumbo GPS con una media m\'ovil
circular y, tras recorrer \texttt{WIND\_OBS\_DISTANCE\_M} $=30$~m, fija el viento
en \texttt{rumbo\_suavizado $+ 90^{\circ}$}. Tambi\'en existe un comando LoRa
\texttt{wind-command} para forzar la direcci\'on manualmente (respaldo fiable).

\begin{warningbox}
  El desfase de $+90^{\circ}$ y el umbral de 30~m necesitan validaci\'on en campo.
  Adem\'as, durante la medici\'on el heartbeat se ve igual que en reposo salvo por
  el campo \texttt{wobs}.
\end{warningbox}

\subsection{Fase 6 --- Algoritmo de navegaci\'on a vela}

El algoritmo de navegaci\'on (\file{navigation.h}) fue desarrollado por un
compa\~nero del grupo (rama \texttt{Testautoboat} del repositorio
\texttt{DeltaGod/PER\_MEA\_LYINCROYAP}) e integrado en este firmware. Implementa:

\begin{itemize}
  \item Zonas prohibidas de ce\~nida y empopada (\emph{upwind/downwind}) con
    zigzag de bordeo.
  \item Evitaci\'on de trasluchada (\emph{empannage}) mediante una m\'aquina de
    estados.
  \item Corredor de tolerancia lateral (\emph{cross-track}) alrededor de la l\'inea
    al waypoint.
  \item Orzar/arribar (\emph{lofer/abattre}) para el seguimiento fino de rumbo.
\end{itemize}

La \'ultima versi\'on integrada aporta: selecci\'on de amura inicial m\'as
inteligente (\texttt{nav\_sideMovingTowardWaypoint}, que elige el lado que m\'as
progresa hacia el waypoint), detecci\'on de llegada por ``plano del waypoint
superado dentro del corredor'' (\texttt{nav\_hasPassedWaypoint}) y un corredor por
defecto ampliado de 20 a 100~m. La envoltura \file{AutoController} mapea el rumbo
de tim\'on del algoritmo 1:1 a grados f\'isicos del winch, limitado a
$\pm20^{\circ}$ (\texttt{ROTOR\_AUTO\_RANGE\_DEG}), ya que el $\pm90^{\circ}$ completo
resultaba demasiado agresivo.

\begin{warningbox}
  La navegaci\'on aut\'onoma est\'a \textbf{codificada pero nunca probada en el
  agua} y no tiene pruebas unitarias. Falta tambi\'en un \emph{WinchTracker}: el
  Regatta ECO II no tiene realimentaci\'on de posici\'on, por lo que la posici\'on
  real del tim\'on no se conoce (ver Secci\'on~\ref{sec:abiertos}).
\end{warningbox}

% =============================================================================
\section{Interfaz de estaci\'on terrena (IHM)}
% =============================================================================

La estaci\'on terrena consta de un segundo T-Beam con \file{transceiver.ino}
(puente USB$\leftrightarrow$LoRa) y una aplicaci\'on web (\file{IHM/}) en Python
FastAPI + MongoDB + un mapa Leaflet. Funciones principales:

\begin{itemize}
  \item Mapa con posici\'on del barco, waypoints (clic para definirlos), env\'io de
    ruta y \emph{polling} de telemetr\'ia a 1~Hz.
  \item Panel de salud con sem\'aforos (Enlace, GPS, Bater\'ia, RC, Control) y RSSI.
  \item Br\'ujula de viento y predicci\'on de trayectoria (r\'eplica en JavaScript
    de \file{navigation.h}).
  \item Barra de progreso de la medici\'on de viento.
  \item Autodetecci\'on del puerto del transceptor por n\'umero de serie del
    CP2104 (transceptor $=$ \texttt{01C00B54}, barco $=$ \texttt{02126CF1}).
\end{itemize}

% =============================================================================
\section{Problemas resueltos y soluciones detalladas}
% =============================================================================

Esta secci\'on re\'une, de forma exhaustiva, los problemas significativos
encontrados durante el desarrollo y c\'omo se resolvieron.

\subsection{Alimentaci\'on del T-Beam --- imposible alimentar desde los pines}
\label{sec:alim}

\begin{warningbox}
  \textbf{Problema:} el T-Beam V1.1 \textbf{no se puede alimentar
  inyectando 5~V por los pines del header de expansi\'on}. Las \'unicas v\'ias de
  alimentaci\'on v\'alidas son el \textbf{portabater\'ias 18650/LiPo} (conector
  JST de la propia placa) y el \textbf{puerto USB}. El PCB V2, en cambio, fue
  dise\~nado para alimentar el T-Beam a trav\'es del pin \texttt{+5V} del header de
  expansi\'on (salida del regulador reductor Pololu).

  El motivo es la topolog\'ia interna de energ\'ia del T-Beam: el pin \texttt{+5V}
  del header est\'a situado en el camino del VBUS de USB, despu\'es del diodo de
  protecci\'on, y \textbf{no entra de forma fiable al AXP192} (el PMIC que habilita
  los rieles de 3,3~V, GPS y LoRa). En consecuencia, alimentando solo por ese pin
  el AXP192 a menudo \textbf{no arranca}, y bajo carga la tensi\'on se desploma
  (se observ\'o el barco corriendo a $\approx$0,81~V, claramente en
  \emph{brownout}), lo que adem\'as provoca la re-enumeraci\'on continua del
  conversor USB-serie (mensajes de kernel \texttt{cp210x ... status -19}).
\end{warningbox}

\begin{okbox}
  \textbf{Soluci\'on r\'apida adoptada:} se cort\'o un cable USB-A~$\to$~USB-C; el
  extremo USB-A (l\'ineas \texttt{+5V} y \texttt{GND}) se \textbf{sold\'o
  directamente al riel de 5~V del PCB} (salida del Pololu), y el extremo USB-C se
  conecta al puerto USB del T-Beam. As\'i, la energ\'ia del PCB entra al T-Beam por
  su \textbf{camino de alimentaci\'on USB v\'alido} (VBUS $\to$ AXP192), en lugar de
  por el pin \texttt{+5V} del header, que no era capaz de arrancar el PMIC.

  \textbf{Recomendaci\'on:} esta soluci\'on es un \emph{parche} funcional pero poco
  elegante (un cable USB soldado al riel). Una correcci\'on definitiva en una
  futura revisi\'on del PCB ser\'ia rutear la alimentaci\'on del PCB directamente a
  la entrada correcta del T-Beam (VBUS o el JST de bater\'ia), evitando el pin
  \texttt{+5V} del header.
\end{okbox}

\subsection{Motor que funciona ``a impulsos'' --- corte por baja tensi\'on}
\label{subsec:lvc}

\begin{warningbox}
  \textbf{S\'intoma:} con la bater\'ia a $\approx$7~V (seg\'un telemetr\'ia) el
  motor no gira de forma continua: arranca un instante, se corta, y solo vuelve a
  funcionar si se lleva el mando a 0~\% y luego de nuevo a 100~\% un tiempo
  despu\'es.

  \textbf{Causa (no es un error de software):} es el \textbf{corte por baja
  tensi\'on (LVC)} del ESC. Una bater\'ia 2S a 7~V est\'a ya baja
  ($\approx$3,5~V/celda); bajo la carga del motor la tensi\'on cae por debajo del
  umbral de corte del ESC ($\approx$3,0--3,2~V/celda) y el ESC corta el motor para
  proteger el pack. Sin carga la tensi\'on se recupera, el corte se rearma y el
  motor vuelve a girar un momento, produciendo justamente el comportamiento ``a
  impulsos''. El efecto empeora mucho si el ESC tiene mal configurado el n\'umero
  de celdas.
\end{warningbox}

\begin{okbox}
  \textbf{Soluci\'on:} encadenar en serie 2 baterias LiPo ($\approx$15,4~V).
  Ning\'un cambio de c\'odigo corrige un corte por baja tensi\'on; el firmware
  env\'ia una consigna estable.
\end{okbox}

\subsection{Fallo al programar el barco --- brownout durante el flash}

\begin{warningbox}
  \textbf{S\'intoma:} la programaci\'on falla repetidamente (``\emph{chip stopped
  responding}'', ``\emph{No serial data received}'') y el CP2104 se re-enumera en
  mitad de la escritura. Esto solo pasa cuando la T-Beam esta conectada directamente al PCB.

  \textbf{Causa:} el mismo \emph{brownout} de alimentaci\'on. El pico de corriente
  del borrado/escritura de flash desploma la tensi\'on y el ESP32 se resetea.
\end{warningbox}

\begin{okbox}Desconectar la T-beam del PCB, Programarla y sin desconectarla del USB volver a colocarla en el PCB.
\end{okbox}

\subsection{Otros problemas de firmware ya resueltos}

\subsubsection{DUTY\_MODE\_1 en OPR\_B --- servo del tim\'on inoperante}

Doble inversi\'on a trav\'es del BC548 (ver Secci\'on~\ref{subsec:bc548}). Soluci\'on:
\texttt{DUTY\_MODE\_0} en ambos canales.

\subsubsection{Umbral de armado de ESC demasiado bajo}

El m\'inimo real del manche de gas del PTR-6A es $\approx$\us{1265}, por lo que el
umbral de armado de \us{1050} nunca se alcanzaba. Se subi\'o a \us{1300}.
\textbf{Nota:} el gesto de armado se \textbf{elimin\'o} despu\'es, al pasar al modo
manual unificado y al control bidireccional (era incompatible con el nuevo mapeo
del gas).

\subsubsection{Modo bloqueado en Automatic --- confusi\'on CH4/CH5}

La versi\'on inicial le\'ia CH4 para el modo, cuando el conmutador est\'a en CH5.
Soluci\'on: leer \texttt{frame.ch5}.

\subsubsection{Puerto serie ocupado al programar}

Un monitor serie o script Python manten\'ia el puerto abierto
(\texttt{Errno 16: Device or resource busy}). Soluci\'on: cerrar el monitor o
\texttt{fuser /dev/ttyUSB0} y matar el proceso.

% =============================================================================
\section{Problemas que encontramos --- opciones de mejora}
\label{sec:abiertos}
% =============================================================================

Esta secci\'on re\'une los problemas a\'un \textbf{abiertos} y, para cada uno, la
opci\'on de mejora propuesta para resolverlo. Son las l\'ineas de trabajo
prioritarias para que el dron pueda cumplir realmente su misi\'on.

\subsection{Alcance de LoRa muy pobre dentro de la caja de electr\'onica}

\begin{warningbox}
  \textbf{Problema:} con la antena LoRa actual ubicada \textbf{dentro de la caja de
  electr\'onica}, el alcance de radio es muy pobre. La caja act\'ua como pantalla y
  la proximidad de la electr\'onica degrada el patr\'on de radiaci\'on, reduciendo
  dr\'asticamente el RSSI y el rango \'util de telemetr\'ia y comandos.
\end{warningbox}

\begin{okbox}
  \textbf{Mejora propuesta:} adquirir una \textbf{antena LoRa externa} (433~MHz) y
  ubicarla \textbf{sobre el m\'astil principal}, lo m\'as alta y despejada posible,
  fuera de la caja. Esto mejora la l\'inea de vista y aleja la antena de las masas
  met\'alicas y del ruido de la electr\'onica.

  \textbf{A considerar --- conector:} hay que verificar el conector de la antena y
  del m\'odulo. El SX1276 del T-Beam suele exponer \textbf{u.FL/IPEX}, mientras que
  las antenas externas de 433~MHz traen habitualmente \textbf{SMA}. Ser\'a
  necesario el \textbf{cable adaptador / pigtail correspondiente (u.FL~$\to$~SMA)} y
  cuidar que sea de 433~MHz (no de 2,4~GHz) para no degradar la adaptaci\'on.
\end{okbox}

\subsection{Inconsistencias del GPS}

\begin{warningbox}
  \textbf{Problema:} el GPS presenta inconsistencias (fix intermitente,
  pocos sat\'elites, posici\'on poco estable), agravadas por la ubicaci\'on de la
  antena dentro o cerca de la caja y la electr\'onica.
\end{warningbox}

\begin{okbox}
  \textbf{Mejora propuesta:} adquirir \textbf{otra antena externa, esta vez de
  GPS}, y ubicarla \textbf{sobre la cubierta}, con vista despejada del cielo y lejos
  de la antena LoRa y de la electr\'onica.

  \textbf{A considerar --- conector:} igual que con LoRa, hay que verificar el
  conector. La entrada de antena GPS del NEO-6M es \textbf{u.FL}; las antenas GPS
  activas suelen venir con \textbf{SMA} o directamente con cable u.FL. Habr\'a que
  prever el \textbf{adaptador adecuado} y confirmar que la antena sea
  \textbf{activa} (el supervisor de antena ya est\'a habilitado por software, ver
  Secci\'on~\ref{subsec:antena}).
\end{okbox}

\subsection{Sin medici\'on de viento ni de la posici\'on de la vela}

\begin{warningbox}
  \textbf{Problema:} el esquema actual \textbf{no mide de d\'onde viene el viento}
  ni \textbf{la posici\'on real de la vela}. Como el winch del tim\'on (Regatta ECO
  II) tampoco tiene realimentaci\'on de posici\'on, \textbf{no es posible conocer la
  posici\'on real del tim\'on / safran} ni cerrar el lazo de control con datos
  reales: toda la navegaci\'on aut\'onoma depende de estimaciones (viento deducido
  del track GPS, posici\'on de tim\'on supuesta a partir de la consigna).
\end{warningbox}

\begin{okbox}
  \textbf{Mejora propuesta:} conseguir \textbf{2 sensores} que midan
  respectivamente la \textbf{direcci\'on del viento} y la \textbf{posici\'on de la
  vela}, y conectarlos al \textbf{bus I2C de la placa} a trav\'es del
  \textbf{conector ya previsto en el PCB} (\gpio{21}/\gpio{22}, que el firmware
  dej\'o libres al eliminar \texttt{Wire.end()}). Con la direcci\'on del viento real
  y la posici\'on real de la vela, s\'i ser\'ia posible deducir la posici\'on
  efectiva del tim\'on y cerrar correctamente el lazo de navegaci\'on.

  \textbf{Ventaja arquitect\'onica:} el firmware ya reserva el bus I2C y el PCB ya
  expone el conector de expansi\'on, por lo que la integraci\'on el\'ectrica es
  directa; faltar\'ia el driver \file{SensorBus} en software.
\end{okbox}

\subsection{Poco espacio para las bater\'ias en la caja}

\begin{warningbox}
  \textbf{Problema:} la caja de electr\'onica tiene \textbf{poco espacio para las
  bater\'ias}, lo que limita la capacidad embarcada y, por tanto, la autonom\'ia y
  el margen de tensi\'on bajo carga (relacionado directamente con el corte por baja
  tensi\'on de la Secci\'on~\ref{subsec:lvc}).
\end{warningbox}

\begin{okbox}
  \textbf{Mejora propuesta:} revisar el dise\~no mec\'anico de la caja para alojar
  un pack de mayor capacidad (o una mejor distribuci\'on de las celdas), teniendo
  en cuenta adem\'as el espacio que ocupar\'an los nuevos sensores y los conectores
  de las antenas externas. Conviene dimensionar el pack para mantener la tensi\'on
  por encima del umbral LVC del ESC durante toda la misi\'on.
\end{okbox}

% =============================================================================
\section{Estado actual del sistema}
% =============================================================================

\begin{longtable}{C{1.4cm} L{5.3cm} C{2.6cm} L{3.6cm}}
  \toprule
  \rowcolor{tableheadbg}
  \textcolor{tableheadfg}{\textbf{Fase}} &
  \textcolor{tableheadfg}{\textbf{Objetivo}} &
  \textcolor{tableheadfg}{\textbf{Estado}} &
  \textcolor{tableheadfg}{\textbf{Notas}} \\
  \midrule
  \endhead
  \bottomrule
  \endfoot
  1 & Control manual RC & \done\ Completo & Modo unificado + motor bidireccional \\[3pt]
  2 & GPS + navegaci\'on & \done\ Completo & Fix probado en exterior \\[3pt]
  3 & LoRa + telemetr\'ia & \done\ Completo & TX/RX verificados \\[3pt]
  PCB V2 & Cableado nuevo + I2C libre & \done\ Completo & Compilado y probado \\[3pt]
  5 & Estimaci\'on de viento & \done\ Codificado & Sin validar en agua \\[3pt]
  6 & Navegaci\'on aut\'onoma & \done\ Codificado & Sin probar en agua, sin tests \\[3pt]
  4 & Sensores I2C & \todo\ Pendiente & Ver Secci\'on~\ref{sec:abiertos} \\[3pt]
  7 & Registro / misi\'on completa & \todo\ Pendiente & SPIFFS, propulsi\'on auto \\
\end{longtable}

\subsection{Valores de calibraci\'on actuales (resumen)}

\begin{longtable}{L{5cm} C{2.4cm} L{5cm}}
  \toprule
  \rowcolor{tableheadbg}
  \textcolor{tableheadfg}{\textbf{Par\'ametro}} &
  \textcolor{tableheadfg}{\textbf{Valor}} &
  \textcolor{tableheadfg}{\textbf{Notas}} \\
  \midrule
  \endhead
  \bottomrule
  \endfoot
  \texttt{SAIL\_CENTER\_US} & \us{1520} & Centro Futaba (no 1500) \\[2pt]
  \texttt{SAIL\_PLUS/MINUS\_US} & 1575 / 1465 & $\pm10^{\circ}$ --- calibrar en banco \\[2pt]
  \texttt{ROTOR\_MIN/MAX\_US} & 1417 / 1583 & $\pm90^{\circ}$ f\'isico (manual) \\[2pt]
  \texttt{ROTOR\_AUTO\_RANGE\_DEG} & $20^{\circ}$ & Recorrido usado en auto \\[2pt]
  \texttt{ESC\_REVERSE\_MIN\_US} & \us{1000} & $-100~\%$ (atr\'as, bidireccional) \\[2pt]
  \texttt{ESC\_NEUTRAL\_US} & \us{1500} & Neutro / paro \\[2pt]
  \texttt{ESC\_MAX\_US} & \us{2000} & $+100~\%$ (adelante) \\[2pt]
  \texttt{CH3\_CENTER\_US} & \us{1545} & Centro del gas --- verificar en banco \\[2pt]
  \texttt{ESC\_SLEW\_US} & 30~\textmu s/tick & Limitaci\'on de variaci\'on \\[2pt]
  Raz\'on divisor bater\'ia & 5,683 & (562k + 120k) / 120k \\
\end{longtable}

% =============================================================================

\end{document}


